\section*{Bab \ref{ch:b2:biogeochemical-cycles} --- Daur Biogeokimia}

\begin{solution}{exo:b2:biogeochemical-cycles:1}
Sebuah \omterm{def:b2:biogeochemical-cycles:reservoir}{tandon} adalah sebuah simpanan unsurnya (sebuah massa); sebuah
\omterm{def:b2:biogeochemical-cycles:reservoir}{fluks} adalah sebuah laju pemindahan antartandon (massa tiap tahun);
\omterm{def:b2:biogeochemical-cycles:reservoir}{waktu tinggalnya} adalah massa \omterm{def:b2:biogeochemical-cycles:reservoir}{tandonnya} dibagi seluruh aliran
keluarnya, yaitu waktu rerata sebuah atom berada di sana. Tumbuhan
darat: $450/120 \approx 4$ tahun (karbon sehelai daun, kira-kira; \omterm{def:b2:plant-meristems:wood}{kayu}
tinggal berpuluh tahun dan reratanya menyembunyikan sebarannya). Laut
dalam: $37000/100 \approx 400$ tahun --- yaitu waktu sebelum sebuah
molekul yang telah tenggelam melihat permukaannya lagi.
\end{solution}

\begin{solution}{exo:b2:biogeochemical-cycles:2}
$2.5\times 2.13 = \qty{5.3}{GtC}$ tiap tahun. $10/2.13 = \qty{4.7}{ppm}$
--- jika semuanya tinggal di udaranya; pada \omterm{prop:b2:biogeochemical-cycles:carbon}{bagian yang mengudara}
\qty{45}{\%}, \qty{2.1}{ppm}.
\end{solution}

\begin{solution}{exo:b2:biogeochemical-cycles:3}
Membuka: fiksasi hayati oleh nitrogenase (sianobakteri, rizobium,
\emph{Azotobacter}, \emph{Frankia}) dan petir. Mengembalikan:
\omterm{def:b2:microbial-metabolism:anaerobic}{denitrifikasi} (anaerob fakultatif seperti \emph{Pseudomonas} dan
\emph{Paracoccus} di dalam tanah dan endapan yang anoksik) dan anamoks
(planktomisetes). Di antaranya: amonifikasi nitrogen organik menjadi
amonium oleh pengurainya (jamur, bakteri); nitrifikasi amonium menjadi
nitrit (\emph{Nitrosomonas}, arkea pengoksidasi amonia) dan nitrit
menjadi nitrat (\emph{Nitrobacter}); serta asimilasi amonium atau
nitrat oleh tumbuhan dan mikroba.
\end{solution}

\begin{solution}{exo:b2:biogeochemical-cycles:4}
Sepanjang waktu geologis, nitrogen selalu dapat dibuat dari
$\mathrm{N_2}$ udaranya yang tak habis-habis oleh pemfiksasinya, yang
disukai setiap kali nitrogennya kurang; fosfor tidak memiliki \omterm{def:b2:biogeochemical-cycles:reservoir}{tandon}
semacam itu dan pasokannya dipancangkan pelapukannya, sehingga
langit-langit keproduktifan jangka panjangnya adalah fosfor. Di dalam
satu musim, fiksasinya lambat dan mahal, dan kebanyakan tumbuhan tidak
dapat mengerjakannya, sehingga nitrogen yang tersedialah yang membatasi
pertumbuhannya di sini dan sekarang --- dan itulah sebabnya pupuknya
kebanyakan nitrogen.
\end{solution}

\begin{solution}{exo:b2:biogeochemical-cycles:5}
$\tau = 1/0.2 = 5$ tahun; $M^{*} = F\tau = \qty{10}{t}$, yakni
$\qty{1}{g/m^3} = \qty{1000}{mg/m^3}$, yang sangat eutrof. Waktu sampai
\qty{90}{\%}: $\tau\ln 10 = 11.5$ tahun. Menyetengahkan masukannya
menyetengahkan \omterm{def:b2:biogeochemical-cycles:reservoir}{keadaan tunaknya}, menjadi \qty{500}{mg/m^3}, yang
dicapai pada skala waktu yang sama.
\end{solution}

\begin{solution}{exo:b2:biogeochemical-cycles:6}
Tahun 1960-an: \qty{0.9}{ppm/yr}, \qty{1.9}{GtC/yr}. Tahun 2010-an:
\qty{2.3}{ppm/yr}, \qty{5.0}{GtC/yr}. Laju pertumbuhannya telah naik
dua setengah kali, seiring dengan pancarannya.
\end{solution}

\begin{solution}{exo:b2:biogeochemical-cycles:7}
\qty{6}{ppm} adalah \qty{13}{GtC}, berbanding produksi primer bersih
\qty{35}{GtC}: gigi gergajinya adalah sepertiganya. Ia lebih kecil
karena respirasi dan penguraiannya berlanjut sepanjang musim panasnya
(gigi gergajinya mencatat penyerapan bersihnya, bukan penyerapan
kasarnya), karena musim belahan bumi selatannya berlawanan lalu
sebagian saling meniadakan, dan karena lautannya meredam ayunannya.
\end{solution}

\begin{solution}{exo:b2:biogeochemical-cycles:8}
$\qty{200}{kg}/(\qty{28}{g/mol}) = 7100$~mol $\mathrm{N_2}$; $16\times
7100 = \num{114000}$~mol ATP; $/30 = 3800$~mol glukosa, \qty{690}{kg}.
Sebuah tanaman yang membangun \qty{10}{t} bahan kering memfiksasi
kira-kira \qtyrange{15}{20}{t} gula setelah respirasinya dihitung,
sehingga fiksasinya berongkos sekitar \qty{4}{\%} fotosintatnya ---
yaitu harga kemerdekaan dari nitrogen tanahnya.
\end{solution}

\begin{solution}{exo:b2:biogeochemical-cycles:9}
Nisbah yang tersedia $\mathrm{N:P} = 32/2.2 = 14.5$, di bawah 16
Redfieldnya: nitratnya habis lebih dahulu, dengan
\qty{0.2}{\micro mol/kg} fosfat tersisa. Karbon yang terfiksasi:
$32\times 106/16 =
\qty{212}{\micro mol/kg}$, yaitu kira-kira
\qty{2.5}{mg} karbon tiap kilogram air.
\end{solution}

\begin{solution}{exo:b2:biogeochemical-cycles:10}
$E^{*} = Q\tau = \qty{500}{GtC}$, yaitu \qty{235}{ppm} di atas 285
praindustrinya: kira-kira \qty{520}{ppm}. Lebih buruk dalam kenyataan
karena lubuknya bukan sebuah lungkang cepat tunggal: laut permukaannya
menjenuh (kemampuannya menyerap $\mathrm{CO_2}$ turun ketika karbonatnya
terpakai), lubuk daratannya melemah bersama pemanasannya, dan
penyingkiran yang sebenarnya memiliki sebuah komponen lambat ---
pencampuran laut dalam dan pelapukan batuannya --- dengan tetapan waktu
seribu sampai seratus ribu tahun, sehingga sebagian kelebihannya tidak
mengendur pada $\tau$ berpuluh tahun mana pun.
\end{solution}

\begin{solution}{exo:b2:biogeochemical-cycles:11}
Yang terukur adalah pertukaran $\mathrm{CO_2}$ atmosfernya dengan
lautan dan biosfernya, bukan peluruhannya: $\mathrm{{}^{14}C}$ bomnya
diencerkan ke dalam lungkang yang jauh lebih besar berupa laut
permukaan serta tumbuhan dan tanahnya, yang lalu mengembalikan karbon
biasa. Paruh waktu 11 tahun kelebihannya (dengan waktu pelipatan
kira-kira 16 tahun) adalah skala waktu pertukaran itu --- yaitu
pergantian karbon atmosfernya terhadap laut permukaan dan daratannya,
yang lebih panjang daripada \omterm{def:b2:biogeochemical-cycles:reservoir}{waktu tinggal} kasarnya yang empat tahun
karena sebagian besar yang pergi langsung kembali.
\end{solution}

\begin{solution}{exo:b2:biogeochemical-cycles:12}
Nitrogen: fiksasi alaminya kira-kira 200~Tg tiap tahun; fiksasi
manusianya (Haber--Bosch 120, tanaman 60, pembakaran 30) kira-kira 210:
\omterm{def:b2:biogeochemical-cycles:reservoir}{fluksnya} melipat dua, dan \omterm{def:b2:biogeochemical-cycles:reservoir}{tandon} yang terpengaruh --- tanah, sungai,
laut pesisir, $\mathrm{N_2O}$ udaranya --- kecil dan cepat, sehingga
perubahannya tampak dalam beberapa dasawarsa, sementara \omterm{def:b2:biogeochemical-cycles:reservoir}{tandon}
$\mathrm{N_2}$-nya tak tersentuh. Karbon: pancaran manusia sebesar
11~GtC tiap tahun hanyalah \qty{5}{\%} pertukaran alami kasarnya yang
kira-kira 200, sehingga \omterm{def:b2:biogeochemical-cycles:reservoir}{fluksnya} hanya tersenggol; tetapi \omterm{def:b2:biogeochemical-cycles:reservoir}{fluks}
alaminya saling meniadakan sedangkan \omterm{def:b2:biogeochemical-cycles:reservoir}{fluks} manusianya tidak, dan ia
telah menaikkan \omterm{def:b2:biogeochemical-cycles:reservoir}{tandon} atmosfernya sebesar \qty{50}{\%} --- senggolan
yang kecil itu memiliki efek kumulatif yang besar pada sebuah \omterm{def:b2:biogeochemical-cycles:reservoir}{tandon}
yang lubuk akhirnya lambat. Daur mana yang ``lebih terusik'' bergantung
pada apakah orang mencacah \omterm{def:b2:biogeochemical-cycles:reservoir}{fluksnya} atau penimbunannya.
\end{solution}

\begin{solution}{pb:b2:biogeochemical-cycles:1}
\textbf{1.} $285\times 2.13 = \qty{607}{GtC}$; $412\times 2.13 =
\qty{878}{GtC}$; kenaikannya \qty{271}{GtC}.
\textbf{2.} $271/700 = 0.39$ (bagiannya sepanjang beberapa dasawarsa
terakhir, dengan pancaran tata guna lahan dihitung, lebih dekat kepada
$0.45$).
\textbf{3.} Lautannya, kira-kira seperempat pancarannya, dan
daratannya (pertumbuhan kembali hutan dan pemupukan $\mathrm{CO_2}$),
kira-kira sepertiganya; sisa \qty{430}{GtC} yang tidak berada di
udaranya.
\textbf{4.} $175/38000 = \qty{0.5}{\%}$ --- dapat diabaikan bagi
seluruh lautannya; tetapi penyerapannya sejauh ini terutama masuk ke
lapisan permukaannya (\qty{900}{GtC}), yang di sana ia merupakan
kenaikan sepersepuluh atau lebih, dan kimia karbonatnya mengubah
kenaikan kecil pada $\mathrm{CO_2}$ terlarutnya menjadi kenaikan yang
lebih besar pada kepekatan ion hidrogennya: pH permukaannya telah turun
sebesar $0.1$, yaitu kenaikan keasaman \qty{30}{\%}.
\textbf{5.} Dengan $Q = Q_0 e^{t/T}$, cobalah $E = Ae^{t/T}$: $A/T = Q_0 -
A/\tau$, sehingga $A = Q_0\tau T/(\tau + T)$ dan \omterm{prop:b2:biogeochemical-cycles:carbon}{bagian yang mengudara}
$\dot E/Q = \tau/(\tau + T)$ bersifat tetap. Pancaran yang tumbuh
\qty{2}{\%} tiap tahun ($T = 50$) dengan $\tau = 40$ memberikan $0.44$:
bagiannya tetap karena lubuk dan sumbernya tumbuh bersama-sama.
\textbf{6.} Dengan $Q$ yang tetap, $E \to Q\tau$ dan $\dot E \to 0$:
\omterm{prop:b2:biogeochemical-cycles:carbon}{bagian yang mengudaranya} turun menuju nol ketika lubuknya menyusul ---
di dalam model satu kotaknya; dalam kenyataan lubuknya menjenuh dan
bagiannya turun lebih lambat.
\textbf{7.} \qty{271}{GtC}.
\textbf{8.} $E^{*} = 10\times 100 = \qty{1000}{GtC}$: $285 +
1000/2.13 = \qty{755}{ppm}$.
\textbf{9.} $E(t) = 1000 + (271 - 1000)e^{-t/100}$; pada $t = 80$,
$E = 1000 - 729\times 0.449 = \qty{672}{GtC}$: \qty{600}{ppm}.
\textbf{10.} $Q = 10(1 - t/50)$: penyelesaian khususnya $E_p = at + b$
dengan $a = -20$, $b = 3000$ (dari $a = 10 - b/\tau$ dan $0 = -0.2 -
a/\tau$); $E = 3000 - 20t + (271 - 3000)e^{-t/100}$. Pada $t = 50$:
$2000 - 2729\times 0.607 = \qty{344}{GtC}$, \qty{447}{ppm}. Lalu $E$
meluruh dengan bebas: pada $t = 80$, $344\,e^{-0.3} = \qty{255}{GtC}$,
\qty{405}{ppm}.
\textbf{11.} \qty{600}{ppm} berbanding \qty{405}{ppm}: skenario yang
kedua mengembalikan udaranya, pada tahun 2100, nyaris ke tempatnya hari
ini.
\textbf{12.} Satu $\tau$ tunggal membuat lubuknya terus bekerja pada
laju yang sama pada sebuah kelebihan yang menyusut; dalam kenyataan
lubuk cepatnya (lapisan tercampurnya, hutan yang tumbuh) sudah
mengambil apa yang dapat mereka ambil dan penyingkiran yang tersisa
adalah lewat pencampuran laut dalam dan pelapukan yang lambat dengan
$\tau$ berabad-abad sampai beribu tahun, sehingga penurunan setelah nol
bersihnya jauh lebih lambat --- kepekatannya tetap kira-kira tetap
selama berabad-abad alih-alih turun.
\textbf{13.} Tanamannya \qty{90}{kg}; yang terdenitrifikasi
\qty{37.5}{kg}; yang tercuci \qty{22.5}{kg} nitrogen tiap hektar tiap
tahun.
\textbf{14.} $\qty{22.5}{kg}/\qty{14}{g/mol} = 1600$~mol N; karbonnya
$1600\times 106/16 = \num{10600}$~mol, \qty{128}{kg}.
\textbf{15.} \num{10600}~mol $\mathrm{O_2}$, \qty{340}{kg}.
\textbf{16.} $\num{10600}/0.25 = \num{42000}$~\unit{m^3}: pencucian
satu hektar dapat mengupas oksigen dari empat puluh ribu meter kubik
air dasar tiap tahun --- dan Mississippi mengaliri seratus juta hektar.
\textbf{17.} $37.5\times 0.01 = \qty{0.375}{kg}$ nitrogen sebagai
$\mathrm{N_2O}$, yakni $0.375\times 44/28 = \qty{0.59}{kg}$
$\mathrm{N_2O}$; $\times 270 = \qty{160}{kg}$ setara $\mathrm{CO_2}$.
\textbf{18.} $150\times 4 = \qty{600}{kg}$ $\mathrm{CO_2}$:
pembuatannya melampaui $\mathrm{N_2O}$ ladangnya empat berbanding satu,
dan keduanya bersama-sama adalah tiga perempat ton tiap hektar tiap
tahun.
\textbf{19.} $\tau = 10$ tahun; $M^{*} = 5\times 10 = \qty{50}{t}$;
kepekatannya $50/(5\times 10^{7}) = \qty{1}{g/m^3} =
\qty{1000}{mg/m^3}$.
\textbf{20.} Tiga puluh kali di atas ambangnya: sangat eutrof.
\textbf{21.} $1/\tau = 0.1 + 0.2 = 0.3$: $\tau = 3.3$ tahun; $M^{*} =
5/0.3 = \qty{16.7}{t}$, \qty{330}{mg/m^3}.
\textbf{22.} $M^{*} = 1/0.3 = \qty{3.3}{t}$, \qty{67}{mg/m^3} --- masih
eutrof; dalam \qty{10}{\%} darinya setelah $\tau\ln 10 = 7.7$ tahun.
\textbf{23.} Endapannya bukanlah sebuah lubuk searah: fosfor yang
tersimpan di dalamnya sepanjang tahun eutrofnya dilepaskan kembali di
bawah keadaan anoksik yang diciptakan eutrofikasinya sendiri (fosfat
yang terikat besi melarut ketika besinya tereduksi), sehingga danaunya
memberi makan dirinya sendiri selama bertahun-tahun setelah masukan
luarnya dipangkas --- yaitu pemuatan dalam, sebuah \omterm{def:b2:biogeochemical-cycles:reservoir}{tandon} kedua dengan
\omterm{def:b2:biogeochemical-cycles:reservoir}{waktu tinggal} yang panjang.
\textbf{24.} $\qty{5000}{kg}/\qty{31}{g/mol} = 1.6\times 10^{5}$~mol P;
karbonnya $\times 106 = 1.7\times 10^{7}$~mol, \qty{205}{t} karbon tiap
tahun --- \qty{20}{g/m^2} di seluruh permukaan danaunya jika danaunya
sedalam \qty{5}{m}, yaitu sebuah ledakan yang berat.
\textbf{25.} \omterm{prop:b2:biogeochemical-cycles:carbon}{Bagian yang mengudara} $0.39$ sepanjang 1850--2020; pada
tahun 2100, \qty{600}{ppm} dengan pancarannya ditahan pada
\qty{10}{GtC/yr} dan \qty{405}{ppm} dengan pancarannya dipangkas
menjadi nol pada tahun 2070 (sebuah model $\tau$ tunggal yang
optimistis); nitrogen yang tercuci \qty{22.5}{kg} tiap hektar tiap
tahun.
\end{solution}
